%; whizzy paragraph -pdf xpdf -latex ./whizzypdfptex.sh
%; whizzy-paragraph "^\\\\begin{frame}\\|\\\\emtext"
% latex beamer presentation.
% platex, latex-beamer でコンパイルすることを想定。 

%     Tokyo Debian Meeting resources
%     Copyright (C) 2012 Junichi Uekawa

%     This program is free software; you can redistribute it and/or modify
%     it under the terms of the GNU General Public License as published by
%     the Free Software Foundation; either version 2 of the License, or
%     (at your option) any later version.

%     This program is distributed in the hope that it will be useful,
%     but WITHOUT ANY WARRANTY; without even the implied warreanty of
%     MERCHANTABILITY or FITNESS FOR A PARTICULAR PURPOSE.  See the
%     GNU General Public License for more details.
%     You should have received a copy of the GNU General Public License
%     along with this program; if not, write to the Free Software
%     Foundation, Inc., 51 Franklin St, Fifth Floor, Boston, MA  02110-1301 USA

\documentclass[cjk,dvipdfmx,12pt]{beamer}
\usetheme{Tokyo}
\usepackage{monthlypresentation}

%  preview (shell-command (concat "evince " (replace-regexp-in-string "tex$" "pdf"(buffer-file-name)) "&")) 
%  presentation (shell-command (concat "xpdf -fullscreen " (replace-regexp-in-string "tex$" "pdf"(buffer-file-name)) "&"))
%  presentation (shell-command (concat "evince " (replace-regexp-in-string "tex$" "pdf"(buffer-file-name)) "&"))

%http://www.naney.org/diki/dk/hyperref.html
%日本語EUC系環境の時
\AtBeginDvi{\special{pdf:tounicode EUC-UCS2}}
%シフトJIS系環境の時
%\AtBeginDvi{\special{pdf:tounicode 90ms-RKSJ-UCS2}}

\newenvironment{commandlinesmall}%
{\VerbatimEnvironment
  \begin{Sbox}\begin{minipage}{1.0\hsize}\begin{fontsize}{8}{8} \begin{BVerbatim}}%
{\end{BVerbatim}\end{fontsize}\end{minipage}\end{Sbox}
  \setlength{\fboxsep}{8pt}
% start on a new paragraph

\vspace{6pt}% skip before
\fcolorbox{dancerdarkblue}{dancerlightblue}{\TheSbox}

\vspace{6pt}% skip after
}
%end of commandlinesmall

\title{東京エリアDebian勉強会}
\subtitle{第142回 2016年8月度}
\author{杉本 典充}
\date{2016年8月20日}
\logo{\includegraphics[width=8cm]{image200607/openlogo-light.eps}}

\begin{document}

\begin{frame}
\titlepage{}
\end{frame}

\begin{frame}{Agenda}
 \begin{minipage}[t]{0.45\hsize}
  \begin{itemize}
  \item 最近あったDebian関連のイベント報告
	\begin{itemize}
	\item 第141回東京エリアDebian勉強会
	\end{itemize}
  \item 事前課題発表
  \end{itemize}
 \end{minipage}
 \begin{minipage}[t]{0.45\hsize}
  \begin{itemize}
   \item Debian Trivia Quiz
   \item 杉本さん「Debianでlxcをセットアップしてみよう」
  \end{itemize}
 \end{minipage}
\end{frame}

\section{イベント報告}
\emtext{イベント報告}

\begin{frame}{第141回東京エリアDebian勉強会 }

\begin{itemize}
\item 場所はサイボウズさんをお借りして開催しました
\item 参加者は12名
\item セミナ内容は、以下内容を発表しました
  \begin{enumerate}
  \item Rogerさん「Debian Installer Screen対応」
  \item yutannihilationさん「go-apt-cacher/mirror」
  \end{enumerate}
\item 残りの時間でhack timeを行い、成果発表をしました
\item \url{https://tokyodebianmeeting.titanpad.com/ep/pad/export/141/rev.1727?format=txt}
\end{itemize} 
\end{frame}

\section{事前課題}
\emtext{事前課題}
{\footnotesize
 \begin{prework}{ mkouhei }
  \begin{enumerate}
  \item $B$"$k(B
  \item LXC
  \item $B%Q%C%1!<%8%a%s%F%J%s%9(B
  \end{enumerate}
\end{prework}

\begin{prework}{ dictoss }
  \begin{enumerate}
  \item $B$"$k(B
  \item lxc
  \item preseed$B!"(Bkfreebsd$B$NF08~D4::(B
  \end{enumerate}
\end{prework}

\begin{prework}{ $B$+$$(B }
  \begin{enumerate}
  \item $B$J$7(B
  \item ($B2sEz$J$7(B)
  \item .deb $B$r$^$@:n$C$?$3$H$,$J$$$N$G:n$C$F$_$k(B
  \end{enumerate}
\end{prework}

\begin{prework}{ kenhys }
  \begin{enumerate}
  \item $B$"$k(B
  \item Docker
  \item $BL$Dj(B
  \end{enumerate}
\end{prework}

\begin{prework}{ koedoyoshida }
  \begin{enumerate}
  \item $B$"$k(B
  \item LXC,Docker$B$*;n$7DxEY!"(Bchroot$B$NJ}$,;H$C$F$k(B
  \item $BL$Dj(B
  \end{enumerate}
\end{prework}

}

\subsection{問題}

%; whizzy-master ../debianmeetingresume201311.tex
% $B0J>e$N@_Dj$r$7$F$$$k$?$a!"$3$N%U%!%$%k$G(B M-x whizzytex $B$9$k$H!"(Bwhizzytex$B$,MxMQ$G$-$^$9!#(B
%

\santaku
{unstable$B$K$F(BGCC$B$,(BTransition$B$5$l$^$7$?!#$5$F!"%P!<%8%g%s$O2?$K(BTransition$B$5$l$?$G$7$g$&$+(B}
{5.4}
{6.1}
{7.0}
{B}
{debian-devel-announce$B$K$F(BGCC-6$B$X(BTransition$B$5$l$^$7$?!#(BGCC-6$B$G$O(BC++14$B$,<BAu$5$l$F$$$^$9!#(BGNU Web$B%5%$%H$N!V(BPorting to GCC 6$B!W$N5-;v$r3NG'$7$F?7(BGCC$B$X0\9T$7$F$/$@$5$$!#%a!<%k!'!V(BTransition news: GCC 6 enabled by default$B!W(B\url{https://lists.debian.org/debian-devel-announce/2016/08/msg00001.html}}

\santaku
{unstalbe$B$K$F(Bperl$B$,(BTransition$B$5$l$^$7$?!#$5$F!"%P!<%8%g%s$O2?$K(BTransition$B$5$l$?$G$7$g$&$+(B}
{5.22}
{5.24}
{6.0}
{B}
{2016$BG/(B5$B7n$K%j%j!<%9$5$l$?(Bperl-5.24$B$X(BTransition$B$5$l$^$7$?!#(B2015$BG/%/%j%9%^%9$K$O(BPerl-6$B%j%j!<%9$NBg%K%e!<%9$,$"$j$^$7$?$,!"(BDebian$B$G:NMQ$5$l$k$K$O@h$,D9$=$&$G$9!#%a!<%k!'!V(Btransition: perl$B!W(B\url{https://bugs.debian.org/cgi-bin/bugreport.cgi?bug=830200}}

\santaku
{Linux Kernel$B$K$*$$$F<!$N(BLTS(=LongTermSupport)$B$H$9$k%P!<%8%g%s$N8uJd$,>e$,$C$F$-$^$7$?!#$5$F%P!<%8%g%s$O2?$G$7$g$&$+(B}
{4.8}
{4.9}
{4.10}
{B}
{stretch$B$N%j%j!<%97W2h$G$O(Blinux-4.10$B$r:NMQ$9$kM=Dj$H$7$?$?$a%U%j!<%:;~4|$rCY$i$;$k$3$H$K$7$F$$$^$7$?!#(BLTS$B$J%+!<%M%k$N%j%j!<%9$,Aa$$$+$i$H$$$C$F$9$G$K%"%J%&%s%9:Q$N%U%j!<%:;~4|$,Aa$^$k$3$H$O$J$$$H;W$$$^$9$,!"Cm;k$7$F$*$$$?$[$&$,$$$$$+$b$7$l$^$;$s!#%a!<%k!'!V(BLinux 4.9 will be next LTS$B!W(B\url{https://lists.debian.org/debian-release/2016/08/msg00147.html}}


\section{Debianでlxcをセットアップしてみよう}
\emtext{Debianでlxcをセットアップしてみよう}

\section{Hack time}
\emtext{Hack time}

\begin{frame}{成果を記入下さい!}

  今回、Hack time時の成果を記録に残してみます。

\url{https://tokyodebianmeeting.titanpad.com/142}

に皆さんアクセス頂き（認証不要です）、各自成果を17:30までに記録するようにお願いします。

\end{frame}
  
\section{今後のイベント}
\emtext{今後のイベント}
\begin{frame}{今後のイベント}
\begin{itemize}
\item 9/17(土) 東京エリアDebian勉強会(予定)
  \begin{itemize}
  \item 発表者募集中
  \end{itemize}
\item 10/15(土) 東京エリアDebian勉強会(予定)
\item 11/5(土),6(日) OSC 2016 Tokyo/Spring
\end{itemize}
\end{frame}

\end{document}

;;; Local Variables: ***
;;; outline-regexp: "\\([ 	]*\\\\\\(documentstyle\\|documentclass\\|emtext\\|section\\|begin{frame}\\)\\*?[ 	]*[[{]\\|[]+\\)" ***
;;; End: ***
